% ============================================================================
% Figure legend for Fig_SomaticIBM_main (svg / pdf / png).
% Final manuscript numbering: Fig. 2. Model definition, update rules and
% scope statements are documented in this repository (README.md and
% data/README.md) rather than in the legend. Format follows the manuscript
% convention ("Fig. 2 | Title." followed by lowercase bold panel letters
% a, b, c, ...).
% ============================================================================

\noindent\textbf{Fig.~2\,|\,The somatic channel of module-integrity
erosion: compartment-network geometry attenuates selection against
core-gene loss.}

\smallskip
\noindent\textbf{a}, Model schematic (illustrative dynamics, not
simulation output). The somatic body is idealised as a $12\times12$
lattice of compartments, each holding $K=8$ nuclei
($N_{\mathrm{tot}}=1{,}152$); three compartments are drawn as rows, and
shading gives the local frequency of nuclei that have irreversibly lost
a bottleneck core gene. Blue dots mark loss-of-function mutation of the
core gene ($\mu=3\times10^{-4}$ per nucleus per nuclear generation);
yellow dots mark nuclear migration through hyphal fusion. Four
architectures are used throughout: C0, fragmented (no fusion); C1, local
fusion between neighbouring compartments (exchange probability
$m_{\mathrm{loc}}=0.08$); C2, local fusion plus rare long-range
anastomosis ($m_{\mathrm{glob}}=0.03$); and C3, strong long-range
fusion ($m_{\mathrm{glob}}=0.25$).

\smallskip
\noindent\textbf{b}, Spatial organisation of core-gene loss at a common
checkpoint ($t=1{,}800$ nuclear generations; $S=0.004$,
$\mu=3\times10^{-4}$, $R=40$ replicates per level), when the
ensemble-mean burden is $\bar{x}=0.31$, $0.42$, $0.42$ and $0.42$ for
C0--C3. Top: compartment grids of the replicate whose whole-body mutant
frequency is closest to the ensemble mean (replicate burdens
$\bar{x}=0.27$--$0.46$); colour bar: mutant fraction per compartment.
Bottom: sectoring coefficient $F=\mathrm{Var}(x_c)/[\bar{x}(1-\bar{x})]$
of every polymorphic replicate at the same checkpoint (dots;
$0.02<\bar{x}<0.98$; $n=19$, $22$, $21$, $20$ of $40$; horizontal bars:
medians). Median sectoring declines monotonically with connectivity:
$F=1.00$, $0.75$, $0.65$, $0.29$ for C0--C3, consistent with the
structural sectoring of the independent neutral calibrations
($F_{\mathrm{struct}}=1.00$, $0.76$, $0.63$, $0.30$); because $F$ is
normalised by $\bar{x}(1-\bar{x})$, the comparison is insensitive to
the residual burden differences between architectures. Dashed line:
$F=1/K=0.125$. The full time course of the grids is shown in
Supplementary Fig.~1.

\smallskip
\noindent\textbf{c}, Fixation probability of a single mutant nucleus
versus the deleterious effect of core-gene loss, $S$ (points:
$R=50{,}000$ replicates per level and value of $S$, conditional on
absorption within $T_{\max}=40{,}000$ generations; error bars: 95\%
Wilson score intervals; downward triangles: no fixation observed,
plotted at the detection limit $2/R=4\times10^{-5}$, which occurs for
C0--C2 at $S=0.05$ and $0.1$ and for C3 at $S=0.1$). Dashed grey line:
neutral expectation, $1/N_{\mathrm{tot}}\approx8.7\times10^{-4}$.
Dotted grey line: classical well-mixed prediction with census size
$N_{\mathrm{tot}}$ and unattenuated selection. Solid lines: rescaled
diffusion reference with effective selection
$s_{\mathrm{eff}}=-\gamma\,F_{\mathrm{eff}}\,S$ ($\gamma=0.25$ per
nuclear generation) and variance-effective size $N_{e,\mathrm{var}}$
($F_{\mathrm{eff}}=0.55,\,0.41,\,0.43,\,0.22$
and $N_{e,\mathrm{var}}=553,\,550,\,508,\,480$ for C0--C3).

\smallskip
\noindent\textbf{d}, Median time to complete core-gene loss under
recurrent mutation ($\mu=3\times10^{-4}$, $R=200$ replicates per point)
for three selective regimes. Costly core ($S=0.02$, $w=0$; blue
circles): $8{,}146$, $6{,}756$, $7{,}470$ and $3{,}919$ generations for
C0--C3, with $22.5\%$ (C0), $6.5\%$ (C1) and $2.5\%$ (C2) of replicates
censored at $T_{\max}=30{,}000$ and none for C3. Short-circuited
selection ($S=0$, $w=0$; grey squares): $2{,}855$, $3{,}436$, $3{,}415$,
$3{,}347$. Selfish saving ($S=0$, $w=0.02$; red triangles): $2{,}905$,
$1{,}293$, $971$, $591$. Dashed line: neutral clock,
$1/\mu\approx3{,}333$ generations.

\smallskip
\noindent\textbf{e}, Paired-arm disconnection experiment ($S=0.02$,
$\mu=3\times10^{-4}$, $R=200$): one eroding lineage under strong mixing
(C3) is split at $t^{*}=751$ generations, when the median replicate
crosses $\bar{x}=0.05$ (dashed line); arm B is then disconnected ($m=0$)
while arm A remains connected. Curves show the ensemble-mean mutant
frequency $\bar{x}$ of each arm and the sectoring coefficient $F$ of arm
B. Median completion time, measured from the start of the experiment:
$4{,}636$ generations for the connected arm ($198/200$ replicates
completed) versus $9{,}228$ for the disconnected arm ($163/200$);
measured from $t^{*}$, $3{,}885$ versus $8{,}477$ generations.
Disconnection roughly doubles the erosion time ($2.0\times$ from the
start, $2.2\times$ from $t^{*}$).